\documentclass[10pt,a4paper]{amsart}
\usepackage[margin=19mm]{geometry}
\usepackage{amsmath,amssymb,mathtools}
\usepackage[scr=rsfs,cal=euler]{mathalfa}
\usepackage{xfrac}
\usepackage[unicode,hidelinks,bookmarks=false]{hyperref}
\newcommand{\ie}{i.e.\ }
\newcommand{\cf}{cf.\ }
\newcommand{\resp}{respectively\ }
\newcommand{\Subsection}{\subsection}
\providecommand{\nobreakdash}{\nobreak\hbox{-}}
\setlength{\emergencystretch}{2em}
\multlinegap0pt
\usepackage{amssymb}
\newtheorem{lemma}{Lemma}
\newcommand\DD{{\mathbb D}}
\def\im{\operatorname{Im}}
\def\re{\operatorname{Re}}
\title{A Gehring-Hayman inequality for~strongly pseudoconvex domains}
\author{{\L}ukasz Kosi\'nski, Nikolai Nikolov, Pascal J. Thomas}
\date{Source: arXiv:2303.04071v2 (CC BY 4.0)}
\begin{document}
\maketitle

\noindent\emph{Source and licence.} A Gehring-Hayman inequality for~strongly pseudoconvex domains. Version arXiv:2303.04071v2 (CC BY 4.0), distributed by arXiv under the Creative Commons Attribution 4.0 International licence, http://creativecommons.org/licenses/by/4.0/, as recorded in the arXiv licence metadata for this exact version and shown on the abstract page https://arxiv.org/abs/2303.04071v2. Full licence notice: \texttt{licenses/arxiv-2303.04071-CC-BY-4.0.txt}.\par\smallskip

\noindent\textit{Editorial scope.} Counted unit: the article's lemma lem:mt (source0.tex lines 738--752). The article's complete original proof of it is delivered with it (source0.tex lines 754--775, one \texttt{proof} environment of 22 lines). No mathematics is written, repaired or re-derived: the statement and the proof are the article's own lines, in the article's own notation, and no retained line is substituted or re-flowed.  No further source-local context is needed: every cross-reference of every retained line resolves inside the two delivered spans.  Nothing was omitted from any retained span, so the package carries no omission record.  No retained line contains a citation, so the wrapper carries no bibliography and no bibliography entry.  No retained line contains a figure environment, an image-inclusion command or drawing code, so no image is delivered and the package declares no asset dependency.  Editorial formatting only, in addition: an AMS \texttt{amsart} preamble that restates the article's own declarations and macros used by the retained lines, namely the environments \texttt{lemma} on separate counters, together with the standard packages those lines need.  The retained environments print in this document as Lemma 1 in order of appearance, whereas the article prints the same items with the numbering of its own full document.  The complete native source of the article as distributed in the arXiv e-print for this exact version is bundled unchanged as \texttt{sources/138307/source0.tex}.
\begin{lemma}\label{lem:mt} 1) Let $\delta:I\to \DD$ be a curve. Let $t\in (0,1)$  be such that $\tilde \delta:=m_t^{-1}\circ \delta $ is contained in $\{\re \lambda \geq \beta\}$, $\beta>-1$. Then 
$$| \delta(s_1) - \delta(s_2)|\simeq (1-t^2) |\tilde \delta (s_1) - \tilde \delta(s_2)|,\quad s_1,s_2\in I.$$ 
Consequently, 
$$l(\delta)\simeq (1-t^2) l(\tilde \delta).$$ 
The estimates above depend only on $\beta$.

2) Suppose additionally
\begin{equation}\label{eq:s}
\left| \tilde\delta(s_2) - \tilde\delta(s_1)\right| \le 
l(\tilde \delta |_{[s_1,s_2]} )\leq (1+ \epsilon) \re(\tilde\delta(s_2) - \tilde\delta(s_1)),
\end{equation} 
where $s_1, s_2\in I$.
%, $s_1<s_2$, are fixed, and that there exists $\gamma>0$  such that $\re(\tilde \delta(s_2) - \tilde \delta(s_1))> \gamma$. 
Assume also that $|\im \tilde \delta(s)|<\epsilon$, $s\in I$,  where $\epsilon>0$ is small enough. Then $$l(\delta|_{[s_1,s_2]}) \leq  C_{\beta} (1+ \epsilon) \re (\delta(s_2) - \delta(s_1)).$$ 
\end{lemma}

\begin{proof}
If $\re x_1, \re y_1\geq -\beta$, then $|1+t x_1|$ and $|1+t y_1|$ are between $1-\beta$ and 2, so by direct computations $|m_t(x_1) -m_t(y_1)|\simeq  (1-t^2)|y_1-x_1|$.

The second part is technical, as well. The first part
and assumption \eqref{eq:s} imply that $l(\delta|_{[s_1, s_2]}) \leq c_\beta (1+\epsilon) (1-t^2) \re (\tilde\delta (s_2) - \tilde \delta (s_1)).$

To get the assertion, compute
\[ 
\re (\delta(s_2) - \delta(s_1)) = (1-t^2) 
\re \left(     
\frac{\tilde\delta(s_2) - \tilde\delta(s_1)}{(1+t\tilde\delta(s_1)) (1+t\tilde\delta(s_2))} \right)
\]
Since $|\im \tilde \delta(s)|<\epsilon$, the argument of the 
denominator is bounded by $C\epsilon$, and by \eqref{eq:s}, so is
the argument of the numerator. So we get
\[
\re (\delta(s_2) - \delta(s_1)) \geq
(1-t^2)(1-C\epsilon)\re \left( \tilde\delta(s_2) - \tilde\delta(s_1)\right),
\]
which finishes the proof.
%If $\epsilon $ is small enough, then $\epsilon\im (\tilde \delta(s_2) - \tilde\delta (s_1))< c \gamma$ for any $c$. Thus $\re (\delta(s_2) - \delta(s_1)) \geq (1-t^2) c'_{\beta}/2 \re( \tilde\delta (s_2) - \tilde \delta (s_1))$.
\end{proof}

\end{document}
